\section{Evaluation}
\label{sec:evaluation}

We evaluate Chiyoda with one primary policy comparison and two supporting
stress studies. The primary study is
\texttt{study\_information\_control.yaml}, a 50-seed comparison of eight
emergency communication regimes: no deliberate intervention, static beacon
broadcast, global PA broadcast, responder relay, entropy-targeted
communication, density-aware communication, exposure-aware communication, and
bottleneck-avoidance communication. The exported bundle contains
\statRunCount{} runs, \statVariantCount{} variant aggregates, and
\statInterventionEvents{} intervention events. Each run uses the same
station-sarin scenario, a 120-agent commuter/tourist population, asymmetric
initial information, gossip, beacon infrastructure, three delayed responders,
and a 500-step horizon.

The two supporting studies each use 30 seeds. The intervention ablation study
varies targeting budget, timing, and bottleneck awareness. The message-quality
study varies broadcast credibility, delay, and repetition. These studies are
not replacements for the main comparison; they test whether the main
information-safety interpretation is robust to natural design choices.

\subsection{Outcome Measures}
\label{subsec:outcomes}

We report three classes of outcome. Physical outcomes include evacuated agents,
mean travel time, incapacitation count, mean hazard exposure, peak bottleneck
queue, and exit imbalance. Information outcomes include mean entropy, entropy
reduction, intervention reach, entropy reduction attributable to interventions,
and belief-accuracy gain. Coupled outcomes include information-safety
efficiency (ISE) and harmful convergence index (HCI). These coupled metrics are
the main object of the paper: they ask whether a policy buys useful belief
improvement without creating exposure or queueing pressure.

Variant tables report means across seeds. Pairwise comparisons against the
no-intervention baseline use two-sided Mann--Whitney tests over run-level
summary rows, with effect sizes used as descriptive support rather than as
standalone proof of operational superiority. This choice is conservative for
the current stage: the simulator is stochastic and stylized, so the evaluation
focuses on robust tradeoff patterns rather than on declaring a universal
winner.

\begin{figure*}[!t]
  \centering
  \includegraphics[width=\textwidth]{../out/information_control_50/figures/13_information_safety_frontier.pdf}
  \caption{Information-safety frontier across communication policies in the
  50-seed main study. Static beacons produce the highest information-safety
  efficiency, while density-aware and exposure-aware policies show the largest
  harmful-convergence penalties under this scenario.}
  \label{fig:safety-frontier}
\end{figure*}

\begin{figure*}[!t]
  \centering
  \includegraphics[width=\textwidth]{../out/information_control_50/figures/12_intervention_timeline.pdf}
  \caption{Intervention reach and entropy effects over time. Broad reach and
  useful safety effect separate: global broadcast reaches many agents, but its
  information-safety efficiency remains below local static beacons, bottleneck
  avoidance, responder relay, and entropy targeting.}
  \label{fig:intervention-timeline}
\end{figure*}

\subsection{Primary Policy Comparison}
\label{subsec:aggregate-findings}

Table~\ref{tab:aggregate} summarizes the 50-seed variant aggregates. The first
result is negative but important: no communication policy dominates the
baseline across physical evacuation outcomes. The no-intervention baseline
evacuates 10.76 agents on average. Static beacons evacuate 10.90, global
broadcast evacuates 11.00, and the remaining active policies range from 9.70
to 10.86. Mann--Whitney tests against the baseline do not show significant
improvements in evacuated count or hazard exposure for any active policy at
$n=50$ seeds. This is why the paper cannot claim that more information
monotonically improves evacuation.

\begin{table*}[!t]
\centering
\caption{Variant aggregates from the 50-seed information-control study. ISE
denotes information-safety efficiency; HCI denotes harmful-convergence index.}
\label{tab:aggregate}
\begin{tabular}{lrrrrrr}
\toprule
Policy & Evac. & Travel (s) & Exposure & Queue & ISE & HCI \\
\midrule
Static beacon & 10.90 & 10.88 & 2.900 & 2.72 & 0.0136 & 7.27 \\
Bottleneck avoidance & 9.70 & 9.76 & 2.904 & 2.84 & 0.0113 & 7.91 \\
Responder relay & 10.52 & 10.69 & 2.905 & 2.76 & 0.0102 & 6.31 \\
Entropy-targeted & 9.90 & 10.31 & 2.873 & 2.86 & 0.0091 & 7.02 \\
Global broadcast & 11.00 & 11.86 & 2.866 & 2.84 & 0.0058 & 8.28 \\
Density-aware & 10.78 & 10.63 & 2.907 & 2.80 & 0.0050 & 9.41 \\
Exposure-aware & 10.86 & 11.65 & 2.916 & 2.86 & 0.0026 & 8.33 \\
No intervention & 10.76 & 11.53 & 2.919 & 2.74 & 0.0000 & 0.00 \\
\bottomrule
\end{tabular}
\end{table*}

The second result is that communication policy still matters, even when
physical outcomes remain noisy. Static beacons are the best policy by
information-safety efficiency (\statBestPolicy{}, \statBestEfficiency{}). They
combine local reach, repeated exposure, and high credibility: enough to reduce
belief uncertainty without synchronizing the full population toward the same
capacity-constrained route. Responder relay has lower efficiency than static
beacons, but also the lowest HCI among the active policies in the main study
(6.31), which supports the idea that delayed, mobile, high-credibility sources
can be conservative safety controllers.

The third result is that wide reach is not the same as useful reach. Global
broadcast slightly improves mean evacuation count relative to the baseline
(11.00 versus 10.76) and has the lowest mean hazard exposure among the main
policies (2.866). However, it ranks fifth by ISE (0.0058) and has a high HCI
(8.28). In this scenario, the global message is useful but inefficient: it
homogenizes route guidance broadly, while local and capacity-aware policies
produce more information benefit per unit of safety pressure.

The fourth result is that adaptive targeting exposes a real control tradeoff.
Entropy targeting has the largest intervention accuracy gain in the main study
(1.76), and both entropy targeting and bottleneck avoidance significantly
reduce mean travel time against no intervention
($p=0.035$ and $p=0.014$, respectively). Yet these same policies do not improve
evacuation count, and bottleneck avoidance has a near-significant reduction in
evacuated agents ($-1.06$, $p=0.061$). The practical interpretation is that
better-informed trajectories can be shorter for agents who move, while still
failing to improve aggregate completion under capacity and hazard pressure.

\subsection{Intervention Ablation}
\label{subsec:ablation-results}

Table~\ref{tab:ablation} reports the 30-seed ablation study. The strongest
finding is budget non-monotonicity. Early low-budget entropy targeting has the
highest ISE in the ablation (0.0253), more than double early high-budget
entropy targeting (0.0092). The high-budget version improves belief accuracy
more strongly (1.89 versus 1.26), but also produces larger queue pressure and
lower information-safety efficiency. This supports the paper's core mechanism:
more certainty is not automatically safer when many agents receive similar
guidance under shared bottleneck constraints.

\begin{table*}[!t]
\centering
\caption{Supporting intervention ablation over 30 seeds.}
\label{tab:ablation}
\begin{tabular}{lrrrrrr}
\toprule
Condition & Evac. & Travel (s) & Exposure & Queue & ISE & HCI \\
\midrule
Entropy early, low budget & 10.30 & 11.74 & 2.809 & 2.93 & 0.0253 & 8.57 \\
Bottleneck late & 11.17 & 11.57 & 2.873 & 2.77 & 0.0143 & 7.69 \\
Bottleneck early & 10.40 & 10.85 & 2.778 & 2.97 & 0.0140 & 7.98 \\
Entropy early, high budget & 10.70 & 11.65 & 2.879 & 3.00 & 0.0092 & 7.08 \\
Entropy late, high budget & 11.47 & 11.80 & 2.866 & 2.80 & 0.0070 & 7.43 \\
No intervention & 11.20 & 11.84 & 2.869 & 2.67 & 0.0000 & 0.00 \\
\bottomrule
\end{tabular}
\end{table*}

The timing ablation also matters. Late entropy targeting recovers evacuation
count relative to early entropy targeting (11.47 versus 10.70 under the
high-budget setting), but its ISE is lower. Bottleneck targeting is less
sensitive to late deployment: the late bottleneck condition has slightly higher
evacuation count and slightly higher ISE than early bottleneck targeting. This
suggests that bottleneck-aware control is a better candidate for delayed
operational response, while entropy targeting needs stricter budget control.

\subsection{Message Quality}
\label{subsec:message-quality-results}

Table~\ref{tab:message-quality} reports the 30-seed message-quality study.
Credibility has the expected direction for information efficiency:
high-credibility global broadcast has the highest ISE among the
broadcast-quality conditions (0.0066), delayed high-credibility broadcast is
close behind (0.0063), medium credibility falls to 0.0034, and low credibility
falls to 0.0026. Repeating low-credibility messages frequently does not repair
this loss; it has the lowest ISE in the study (0.0009) while preserving high
harmful convergence (8.92).

\begin{table*}[!t]
\centering
\caption{Supporting message-quality study over 30 seeds.}
\label{tab:message-quality}
\begin{tabular}{lrrrrrr}
\toprule
Condition & Evac. & Travel (s) & Exposure & Queue & ISE & HCI \\
\midrule
High-credibility global & 11.13 & 12.01 & 2.889 & 2.83 & 0.0066 & 8.42 \\
Delayed high credibility & 10.80 & 10.70 & 2.856 & 2.80 & 0.0063 & 8.09 \\
Medium-credibility global & 11.53 & 10.85 & 2.862 & 2.80 & 0.0034 & 7.12 \\
Low-credibility global & 12.10 & 11.24 & 2.893 & 2.90 & 0.0026 & 9.07 \\
Frequent low credibility & 11.93 & 11.79 & 2.858 & 2.80 & 0.0009 & 8.92 \\
No intervention & 11.20 & 11.84 & 2.869 & 2.67 & 0.0000 & 0.00 \\
\bottomrule
\end{tabular}
\end{table*}

This support study strengthens the interpretation of global broadcast. Global
messages can increase completed evacuations in some parameterizations, but the
information-safety metric penalizes conditions where the same broad signal also
creates route synchronization. Message quality therefore behaves like a control
parameter, not a cosmetic parameter: credibility and timing change whether
communication improves useful beliefs or merely coordinates crowd movement.

\subsection{Hypotheses Revisited}
\label{subsec:hypotheses-results}

H1 predicted that adaptive targeting would improve information-safety
efficiency over global broadcast. The 50-seed result supports this for static
beacons, bottleneck avoidance, responder relay, and entropy targeting, all of
which exceed global broadcast by ISE. The result is weaker for density-aware
and exposure-aware targeting, which have nonzero efficiency but lower safety
efficiency than global broadcast in this scenario.

H2 predicted that global broadcast would reduce entropy quickly but risk route
synchronization. The data support the synchronization concern more strongly
than the entropy-reduction claim. Global broadcast has high reach and high HCI,
but its ISE trails more localized policies. The message-quality study further
shows that frequent low-credibility broadcast is especially inefficient.

H3 predicted that exposure-aware targeting would reduce hazard exposure.
Exposure-aware targeting does not meaningfully reduce exposure in the 50-seed
main study: its mean hazard exposure is 2.916, close to the baseline value of
2.919 and above several other policies. This hypothesis is not supported by
the present configuration.

H4 predicted that responder relay would be robust but delayed. The result is
partially supported. Responder relay is third by ISE, has the lowest HCI among
the main active policies, and avoids the extreme synchronization penalties of
some adaptive policies. Its limitation is not obvious physical superiority, but
conservative information control.

\subsection{Interpretation}
\label{subsec:interpretation}

The study supports the paper's central framing rather than a single winning
policy. Information interventions should be evaluated as safety-control
actions, not as monotonic entropy minimizers. A policy can reach many agents
and still have low safety efficiency; another can improve belief accuracy and
still fail to improve aggregate completion. The promising direction is
therefore not ``broadcast more,'' but ``broadcast credibly, locally, and with
awareness of capacity and exposure.'' Static beacon and responder relay
policies currently look like conservative baselines, while entropy targeting is
scientifically valuable because it exposes the failure mode created by
overconfident, synchronized route choice.

\subsection{Threats to Evaluation}
\label{subsec:evaluation-threats}

The current evidence is stronger than the pilot run, but still not final
external validation. The experiments use one station geometry, one hazard
scenario, one population size, and a fixed evacuation horizon. The supporting
studies vary intervention design but not architectural layout, sensor latency,
or calibrated pedestrian demographics. The HCI and ISE metrics are also model
constructs; future work should test alternative normalizations and compare
them against trajectory-level bottleneck events, drill data, or expert-coded
evacuation plans.
